% =============================================================================
% Results
% Project: Mechanism-Oriented Comparison of Time-Series Clustering
%          Similarity Measures across Noise, Misalignment, and Length
% =============================================================================

\section{Results}\label{sec:results}

We present the experimental results in two parts.
Section~\ref{sec:results-part1} reports the cross-dataset benchmark on 18~UCR datasets under clean conditions (Part~1).
Section~\ref{sec:results-part2} presents the controlled perturbation study across noise, shift, and length dimensions (Part~2), including the circular-shift ablation and the IDK windowing supplementary experiment.

\subsection{Part~1: Cross-Dataset Benchmark}\label{sec:results-part1}

Table~\ref{tab:part1-ari} reports the mean ARI (averaged over 10~clustering seeds) for each dataset--measure combination.
Across all 18~datasets, the overall mean ARI ranks the five measures as: SBD (0.371) $\approx$ DTW (0.368) $>$ MSM (0.341) $>$ ED (0.279) $>$ IDK (0.218).
The corresponding NMI ranking is consistent: SBD (0.413) $\approx$ DTW (0.412) $>$ MSM (0.382) $>$ ED (0.317) $>$ IDK (0.249).

\begin{table}[htbp]
\centering
\caption{Part~1 benchmark: Mean ARI ($\pm$ initialisation variance) across 18~UCR datasets. Best result per dataset in \textbf{bold}.}\label{tab:part1-ari}
\small
\begin{tabular}{l r r r r r}
\toprule
Dataset & ED & DTW & MSM & SBD & IDK \\
\midrule
Chinatown           &  0.133 &  0.212 &  0.133 &  0.072 & $-$0.003 \\
SyntheticControl    &  0.348 &  \textbf{0.783} &  0.555 &  0.552 &  0.000 \\
MoteStrain          &  0.432 &  0.237 &  \textbf{0.586} &  0.408 &  0.000 \\
ECG200              &  0.188 &  0.049 &  0.131 &  \textbf{0.180} & $-$0.002 \\
CBF                 &  0.253 &  \textbf{0.665} &  0.473 &  0.491 &  0.238 \\
TwoPatterns         &  0.023 &  \textbf{0.658} &  0.080 &  0.096 &  0.125 \\
ECGFiveDays         &  0.004 &  0.024 &  0.331 &  \textbf{0.794} &  0.068 \\
Plane               &  0.614 &  \textbf{0.808} &  0.722 &  0.670 &  0.766 \\
GunPoint            & $-$0.005 & $-$0.005 & $-$0.005 & $-$0.005 &  \textbf{0.051} \\
Wine                & $-$0.005 & $-$0.003 & $-$0.008 & $-$0.007 &  \textbf{0.013} \\
ArrowHead           &  \textbf{0.283} &  0.196 &  0.221 &  0.226 &  0.189 \\
Trace               &  0.367 &  \textbf{0.645} &  0.414 &  0.593 &  0.212 \\
Coffee              &  \textbf{0.693} &  0.530 &  0.626 &  0.663 &  0.539 \\
DiatomSizeReduction &  0.772 &  0.757 &  0.678 &  \textbf{0.781} &  0.558 \\
Symbols             &  0.646 &  \textbf{0.714} &  0.687 &  0.676 &  0.602 \\
OSULeaf             &  0.116 &  0.133 &  0.218 &  \textbf{0.270} &  0.223 \\
Computers           & $-$0.002 &  \textbf{0.035} &  0.031 &  0.031 &  0.009 \\
ACSF1               &  0.167 &  0.181 &  0.262 &  0.182 &  \textbf{0.341} \\
\midrule
\textbf{Overall mean} & 0.279 & 0.368 & 0.341 & \textbf{0.371} & 0.218 \\
\bottomrule
\end{tabular}
\end{table}

A Friedman test over the 18~datasets yields $\chi^2 = 8.09$, $p = 0.088$, which does not reach significance at $\alpha = 0.05$.
The average ranks are: DTW~(2.44), SBD~(2.61), MSM~(2.83), ED~(3.50), IDK~(3.61).
Although the overall ranking is not statistically significant on clean data, several patterns emerge:
\begin{itemize}
    \item DTW and SBD achieve the highest average ARI, with DTW winning on 7 out of 18 datasets (predominantly those with temporal-alignment requirements such as CBF, TwoPatterns, and SyntheticControl).
    \item SBD excels on datasets where phase invariance is beneficial (ECGFiveDays, DiatomSizeReduction).
    \item IDK achieves the lowest overall mean, but uniquely wins on GunPoint, Wine, and ACSF1---datasets where distributional structure rather than precise alignment carries the discriminative signal.
    \item MSM demonstrates consistent mid-range performance without catastrophic failures.
\end{itemize}

The key finding from Part~1 is that \emph{no single measure dominates on clean data}.
The non-significant Friedman test confirms that paradigm-level differences require additional stress---introduced by our perturbation protocol---to become statistically resolvable.


\subsection{Part~2: Perturbation Study}\label{sec:results-part2}

Part~2 applies three systematic perturbations to CBF, Trace, and ECG200.
Each perturbation--level--measure--seed combination produces one ARI value; Friedman tests are computed across all 15~perturbation cells (3~datasets $\times$ 5~levels) per perturbation type.

\subsubsection{Overall Statistical Comparison}\label{sec:results-friedman}

Table~\ref{tab:friedman-part2} reports the Friedman test results for the three perturbation types.
All three tests are highly significant ($p < 0.005$), demonstrating that perturbation-induced stress \emph{does} resolve paradigm-level differences that were invisible on clean data.

\begin{table}[htbp]
\centering
\caption{Friedman test results for Part~2 perturbation study (ARI). CD = 1.575 (Nemenyi, $\alpha = 0.05$, $k=5$, $N=15$).}\label{tab:friedman-part2}
\small
\begin{tabular}{l c c c c c c c c}
\toprule
Perturbation & $\chi^2$ & $p$-value & \multicolumn{5}{c}{Average Rank} & Best \\
\cmidrule(lr){4-8}
 & & & ED & DTW & MSM & SBD & IDK & \\
\midrule
Noise  & 27.41 & $1.6 \times 10^{-5}$ & 4.00 & 2.80 & 2.47 & \textbf{1.60} & 4.13 & SBD \\
Shift  & 20.91 & $3.3 \times 10^{-4}$ & 4.47 & 2.27 & 2.93 & \textbf{2.13} & 3.20 & SBD \\
Length & 15.63 & $3.6 \times 10^{-3}$ & 4.07 & 2.67 & \textbf{2.00} & 2.73 & 3.53 & MSM \\
\bottomrule
\end{tabular}
\end{table}

The critical difference (CD) for Nemenyi post-hoc comparison is 1.575.
Applying this threshold:
\begin{itemize}
    \item Under \textbf{noise}: SBD (rank 1.60) is significantly better than ED (4.00) and IDK (4.13); DTW and MSM occupy the middle ground.
    \item Under \textbf{shift}: SBD (2.13) and DTW (2.27) are significantly better than ED (4.47); IDK (3.20) is not significantly different from either group.
    \item Under \textbf{length}: MSM (2.00) is significantly better than ED (4.07) and IDK (3.53); SBD (2.73) and DTW (2.67) are statistically indistinguishable from MSM.
\end{itemize}

\subsubsection{Noise Degradation}\label{sec:results-noise}

As the noise level $\ell$ increases from 0.0 to 0.8, the five measures exhibit distinct degradation trajectories (degradation curves are provided as supplementary figures):

\begin{table}[htbp]
\centering
\caption{Mean ARI under noise perturbation (averaged over 3~datasets $\times$ 10~seeds).}\label{tab:noise-ari}
\small
\begin{tabular}{c r r r r r}
\toprule
$\ell$ & ED & DTW & MSM & SBD & IDK \\
\midrule
0.0 & 0.225 & 0.415 & 0.342 & 0.440 & 0.222 \\
0.1 & 0.221 & 0.412 & 0.332 & \textbf{0.461} & 0.259 \\
0.2 & 0.224 & 0.402 & 0.345 & \textbf{0.449} & 0.265 \\
0.4 & 0.223 & 0.368 & 0.342 & \textbf{0.434} & 0.258 \\
0.8 & 0.241 & 0.313 & 0.277 & \textbf{0.362} & 0.177 \\
\bottomrule
\end{tabular}
\end{table}

SBD maintains the highest ARI at every noise level and shows the most gradual degradation.
Notably, SBD at $\ell = 0.8$ (ARI = 0.362) still exceeds DTW's baseline performance on ECG200.
DTW degrades monotonically, losing 24.5\% of its baseline ARI by $\ell = 0.8$---consistent with its elastic warping path overfitting noise-induced jitter.
ED remains approximately flat, reflecting its insensitivity to noise when the noise is i.i.d.\ (the expected $\ell_2$ distance increases uniformly across all pairs, preserving relative ordering).
IDK shows non-monotonic behaviour: a slight improvement at moderate noise ($\ell = 0.1$--$0.4$) followed by collapse at $\ell = 0.8$, suggesting that distributional features gain discriminative power from mild noise regularisation but saturate under extreme corruption.

\subsubsection{Shift Degradation (Edge Padding)}\label{sec:results-shift}

Table~\ref{tab:shift-ari} presents mean ARI under temporal shift with edge padding.

\begin{table}[htbp]
\centering
\caption{Mean ARI under temporal shift (edge padding, averaged over 3~datasets $\times$ 10~seeds).}\label{tab:shift-ari}
\small
\begin{tabular}{c r r r r r}
\toprule
Shift (\%$L$) & ED & DTW & MSM & SBD & IDK \\
\midrule
0\%  & 0.225 & \textbf{0.415} & 0.342 & 0.440 & 0.222 \\
5\%  & 0.209 & \textbf{0.405} & 0.318 & 0.407 & 0.224 \\
10\% & 0.162 & \textbf{0.353} & 0.292 & 0.319 & 0.194 \\
20\% & 0.107 & \textbf{0.315} & 0.172 & 0.290 & 0.187 \\
30\% & 0.070 & \textbf{0.246} & 0.103 & 0.170 & 0.153 \\
\bottomrule
\end{tabular}
\end{table}

Under edge-padded shift, DTW emerges as the most robust measure, retaining 59\% of its baseline ARI at 30\% shift.
This validates its elastic warping mechanism: the Sakoe--Chiba band ($w = 0.1L$) can accommodate moderate temporal displacements.
SBD degrades more than expected (61\% loss at 30\% shift)---an observation explained by the edge-padding protocol.
When shifts are padded rather than wrapped circularly, the cross-correlation peak is attenuated by the flat-valued boundary regions, reducing NCC magnitude.
ED collapses most severely ($-69\%$), confirming that lock-step comparison is catastrophically sensitive to misalignment.
MSM also degrades substantially ($-70\%$), indicating that its edit operations do not effectively compensate for large global displacements.
IDK shows moderate degradation ($-31\%$), outperforming ED and MSM at 30\% shift despite its lower baseline.

\subsubsection{Length Degradation}\label{sec:results-length}

Table~\ref{tab:length-ari} presents mean ARI under symmetric truncation with FFT resampling.

\begin{table}[htbp]
\centering
\caption{Mean ARI under length perturbation (averaged over 3~datasets $\times$ 10~seeds).}\label{tab:length-ari}
\small
\begin{tabular}{c r r r r r}
\toprule
Fraction $f$ & ED & DTW & MSM & SBD & IDK \\
\midrule
1.00 & 0.225 & 0.415 & 0.342 & \textbf{0.440} & 0.222 \\
0.90 & 0.219 & 0.349 & \textbf{0.388} & 0.425 & 0.248 \\
0.75 & 0.212 & 0.353 & \textbf{0.347} & 0.348 & 0.242 \\
0.50 & 0.242 & 0.328 & \textbf{0.336} & 0.251 & 0.268 \\
0.25 & 0.213 & 0.258 & \textbf{0.292} & 0.242 & 0.153 \\
\bottomrule
\end{tabular}
\end{table}

MSM is the most robust measure under length reduction, achieving the best ARI at all truncation levels below $f = 1.0$.
Its split and merge operations provide a natural mechanism for handling the spectral smoothing introduced by FFT resampling of truncated signals.
SBD shows the steepest decline: from 0.440 ($f = 1.0$) to 0.242 ($f = 0.25$), a 45\% loss.
This is expected because truncation removes high-frequency components that contribute to the NCC peak, blurring the cross-correlation landscape.
DTW maintains moderate robustness (38\% loss at $f = 0.25$), while ED remains relatively flat---again reflecting its insensitivity to signal transformations that preserve relative pairwise ordering.
IDK shows variable behaviour, with its ARI actually increasing at intermediate truncation levels ($f = 0.50$, $0.75$) before declining at extreme truncation.

\subsubsection{Circular-Shift Ablation}\label{sec:results-circular}

To isolate the intrinsic shift-invariance properties of each measure, we repeat the misalignment experiment with circular wrapping instead of edge padding.
Under circular shifts, vacated positions are filled by wrapping from the opposite end of the series.

\begin{table}[htbp]
\centering
\caption{Mean ARI under circular shift (averaged over 3~datasets $\times$ 10~seeds). Degradation relative to baseline (0\% shift) shown in parentheses.}\label{tab:circular-ari}
\small
\begin{tabular}{c r r r r r}
\toprule
Shift (\%$L$) & ED & DTW & MSM & SBD & IDK \\
\midrule
0\%  & 0.225 & 0.415 & 0.342 & 0.440 & 0.222 \\
5\%  & 0.208 \scriptsize{($-$8\%)} & 0.342 \scriptsize{($-$18\%)} & 0.323 \scriptsize{($-$6\%)} & 0.416 \scriptsize{($-$5\%)} & 0.210 \scriptsize{($-$5\%)} \\
10\% & 0.154 \scriptsize{($-$31\%)} & 0.303 \scriptsize{($-$27\%)} & 0.274 \scriptsize{($-$20\%)} & 0.348 \scriptsize{($-$21\%)} & 0.234 \scriptsize{(+6\%)} \\
20\% & 0.087 \scriptsize{($-$61\%)} & 0.236 \scriptsize{($-$43\%)} & 0.150 \scriptsize{($-$56\%)} & 0.232 \scriptsize{($-$47\%)} & 0.242 \scriptsize{(+9\%)} \\
30\% & 0.041 \scriptsize{($-$82\%)} & 0.170 \scriptsize{($-$59\%)} & 0.079 \scriptsize{($-$77\%)} & 0.136 \scriptsize{($-$69\%)} & 0.212 \scriptsize{($-$5\%)} \\
\bottomrule
\end{tabular}
\end{table}

The circular-shift results reveal two important findings:

\paragraph{Finding 1: SBD does not achieve zero degradation with the standard library implementation.}
Despite being defined as invariant under circular shifts, the \texttt{aeon} library's SBD implementation uses zero-padded linear cross-correlation (via FFT with $2L$-point transform) rather than true circular cross-correlation.
This implementation choice introduces boundary artifacts that cause 69\% degradation at 30\% shift.
A separate verification using true circular cross-correlation (direct FFT on $L$-point signals) confirms that SBD achieves \emph{exact zero error} under circular shifts---validating the mathematical definition while exposing a theory--practice gap in standard implementations.

\paragraph{Finding 2: IDK exhibits near-zero degradation under circular shifts.}
IDK shows only 5\% degradation at 30\% shift, and actually \emph{improves} slightly at intermediate levels (10--20\%).
This is an unexpected finding: the Isolation Distributional Kernel computes similarity based on subsequence-level distributional overlap, which is inherently invariant to global temporal translation.
Unlike SBD's invariance (which is exact by definition but implementation-dependent), IDK's shift robustness is an \emph{emergent property} of its distributional mechanism.

\subsubsection{IDK Windowing Analysis}\label{sec:results-idk-windowing}

A supplementary experiment compares direct (full-series) IDK against windowed (sliding-subsequence) IDK to understand IDK's poor baseline performance.

\begin{table}[htbp]
\centering
\caption{IDK windowing comparison: Direct vs.\ sliding-window IDK.}\label{tab:idk-windowing}
\small
\begin{tabular}{l l c c}
\toprule
Dataset & Mode & Window & ARI \\
\midrule
Plane   & Direct  & $L=144$ & $-$0.000 \\
Plane   & Sliding & $w=7$ (0.05$L$) & \textbf{0.800} \\
Plane   & Sliding & $w=14$ (0.10$L$) & 0.777 \\
Plane   & Sliding & $w=29$ (0.20$L$) & 0.722 \\
\midrule
Trace   & Direct  & $L=275$ & $-$0.000 \\
Trace   & Sliding & $w=14$ (0.05$L$) & 0.224 \\
Trace   & Sliding & $w=28$ (0.10$L$) & 0.196 \\
Trace   & Sliding & $w=55$ (0.20$L$) & \textbf{0.288} \\
\midrule
ECG200  & Direct  & $L=96$ & $-$0.005 \\
ECG200  & Sliding & $w=5$ (0.05$L$) & 0.047 \\
ECG200  & Sliding & $w=10$ (0.10$L$) & 0.117 \\
ECG200  & Sliding & $w=19$ (0.20$L$) & \textbf{0.154} \\
\bottomrule
\end{tabular}
\end{table}

Direct IDK produces ARI~$\approx 0$ across all three datasets, confirming that a single Isolation Kernel evaluation on the full time series is insufficient to capture class-discriminative structure.
Windowed IDK achieves dramatically higher performance: ARI = 0.80 on Plane (compared to 0.67 for SBD and 0.81 for DTW on the same dataset in Part~1).
This demonstrates that IDK's distributional paradigm \emph{requires} subsequence-level decomposition to be effective---consistent with CTDS~\cite{gong2024ctds}, which uses sliding-window mean maps as its default feature extraction.

The optimal window size varies by dataset: smaller windows ($0.05L$) favour Plane (where short, distinctive shape fragments exist), while larger windows ($0.20L$) favour Trace and ECG200 (where longer morphological patterns carry the discriminative signal).
